% !PRINTED 499
\weberpage{499}
Für den Fall $n = 11$ ist die entsprechende Rechnung noch
nicht durchgeführt. Verhältnismäßig einfach erhält man, wenn
man mit Benutzung der Betrachtungsweise des § 119 die Werte
von $\omega$ aufsucht, für die eine der Größen $w_\nu$ verschwindet:
\[
\omega = i, \quad \omega = \sqrt{-7}, \quad \omega = \frac{1+\sqrt{-7}}{1-\sqrt{-7}},
\]
\[
u = \sqrt[4]{2}, \quad u = \sqrt{2}, \quad u = 1,
\]
und daraus
\[
A_{11} = \left(u^{12} - \frac{2^6}{u^{12}}\right)^3\left(u^{12}-\frac{2^{12}}{u^{12}}\right)\left(u^{12}-\frac{1}{u^{12}}\right).
\]

\begin{center}\rule{3cm}{0.4pt}\end{center}
